\SourcePage{495}{498}
\setcounter{section}{102}
\section{Quantization of a Top's Rotation}

The study of the rotational levels of a polyatomic molecule is often complicated by the need to consider rotation together with vibration. As a preliminary problem, we shall examine the molecule's rotation as that of a rigid body, that is, with its atoms ``rigidly fixed'' (a \emph{top}).

Let $\xi\eta\zeta$ be a coordinate system whose axes point along the three principal axes of inertia of the top and rotate with it. The corresponding Hamiltonian is obtained by replacing the components $J_\xi,J_\eta,J_\zeta$ of its angular momentum in the classical expression for the energy by the corresponding operators:
\begin{equation}
 \hat H=\frac{\hbar^2}{2}\left(\frac{\hat J_\xi^2}{I_A}+\frac{\hat J_\eta^2}{I_B}+\frac{\hat J_\zeta^2}{I_C}\right), \tag{103.1}
\end{equation}
where $I_A,I_B,I_C$ are the principal moments of inertia of the top.

The commutation rules for the operators $\hat J_\xi,\hat J_\eta,\hat J_\zeta$, the components of the angular momentum in the rotating coordinate system, are not evident, since the usual derivation of the commutation rule concerns the components $\hat J_x,\hat J_y,\hat J_z$ in a fixed coordinate system. They are nevertheless readily obtained from the formula
\begin{equation}
 (\hat{\mathbf J}\mathbf a)(\hat{\mathbf J}\mathbf b)-(\hat{\mathbf J}\mathbf b)(\hat{\mathbf J}\mathbf a)=-i\hat{\mathbf J}(\mathbf a\times\mathbf b), \tag{103.2}
\end{equation}
where $\mathbf a,\mathbf b$ are any two vectors characterizing the given body (and commuting with each other). This formula is easily

\SourcePage{496}{499}
verified by evaluating the left-hand side in the fixed coordinate system $xyz$, using the general commutation rules between the components of angular momentum and between them and the components of an arbitrary vector.

Let $\mathbf a$ and $\mathbf b$ be unit vectors along the $\xi$ and $\eta$ axes. Then $\mathbf a\times\mathbf b$ is the unit vector along the $\zeta$ axis, and (103.2) gives
\begin{equation}
 \hat J_\xi\hat J_\eta-\hat J_\eta\hat J_\xi=-i\hat J_\zeta. \tag{103.3}
\end{equation}
The other two relations follow in the same way.

Thus, the commutation rules for the operators representing the angular-momentum components in the rotating coordinate system differ from those in the fixed system only in the sign on the right-hand side\footnote{This circumstance expresses the fact that, in their action on the wave function of the top, a rotation of the $xyz$ system is equivalent to an opposite rotation of the $\xi\eta\zeta$ system.}. It follows that all results previously derived from the commutation rules for eigenvalues and matrix elements hold also for $J_\xi,J_\eta,J_\zeta$, with the sole difference that every expression must be replaced by its complex conjugate. In particular, the eigenvalues of $J_\zeta$ (which in this section will be denoted by $k$, to distinguish them from the eigenvalues $J_z=M$) run through $k=-J,\ldots,+J$, where $J$ (an integer!) is the magnitude of the top's angular momentum.

\textbf{Spherical top.} The energy eigenvalues of a rotating top are most easily found when all three principal moments of inertia are equal: $I_A=I_B=I_C=I$. For a molecule this occurs when it has the symmetry of one of the cubic point groups. The Hamiltonian (103.1) becomes
\[
 \hat H=\frac{\hbar^2}{2I}\hat{\mathbf J}^{2},
\]
and its eigenvalues are
\begin{equation}
 E=\frac{\hbar^2}{2I}J(J+1). \tag{103.4}
\end{equation}
Each such energy level is degenerate with respect to the $2J+1$ directions of the angular momentum relative to the top itself (that is, with respect to the values $J_\zeta=k$)\footnote{Here and below we disregard the physically inessential $(2J+1)$-fold degeneracy, always present, with respect to the directions of the angular momentum relative to the fixed coordinate system. When it is included, the total degeneracy of the spherical-top energy levels is $(2J+1)^2$.}.

\SourcePage{497}{500}
\textbf{Symmetric top.} The energy levels are also straightforward to calculate when only two moments of inertia of the top coincide: $I_A=I_B\ne I_C$. This is the case for molecules possessing a symmetry axis of order greater than two. The Hamiltonian (103.1) takes the form
\begin{equation}
 \hat H=\frac{\hbar^2}{2I_A}(\hat J_\xi^2+\hat J_\eta^2)+\frac{\hbar^2}{2I_C}\hat J_\zeta^2
 =\frac{\hbar^2}{2I_A}\hat{\mathbf J}^{2}+\frac{\hbar^2}{2}\left(\frac1{I_C}-\frac1{I_A}\right)\hat J_\zeta^2. \tag{103.5}
\end{equation}
It follows that the energy in a state with definite $J$ and $k$ is
\begin{equation}
 E=\frac{\hbar^2}{2I_A}J(J+1)+\frac{\hbar^2}{2}\left(\frac1{I_C}-\frac1{I_A}\right)k^2, \tag{103.6}
\end{equation}
which determines the energy levels of the symmetric top.

The degeneracy in $k$ that was present for the spherical top is now partly removed. Energies coincide only for values of $k$ that differ in sign alone; these correspond to opposite directions of the angular momentum relative to the top's axis. Hence the energy levels of the symmetric top with $k\ne0$ are twofold degenerate.

The stationary states of a symmetric top are therefore specified by three quantum numbers: the angular momentum $J$, its projection on the top's axis ($J_\zeta=k$), and its projection on the spatially fixed $z$ axis ($J_z=M$); the energy of the top is independent of the last number. In this connection, note that the magnitude of the angular momentum and its projections on a spatially fixed axis and on an axis rigidly attached to the physical system are simultaneously measurable\footnote{These must not be confused with the projections (which cannot be measured simultaneously) on two spatially fixed axes!}. Indeed, the operators $\hat{\mathbf J}^{2}$ and $\hat J_z$ commute not only with one another but also with $\hat J_\zeta=\hat{\mathbf J}\mathbf n$ ($\mathbf n$ is a unit vector along the $\zeta$ axis). This fact can readily be checked by direct calculation, but is also evident in advance. The angular-momentum operator reduces to the operator of an infinitesimal rotation, while the scalar product $\mathbf J\mathbf n$ of two vectors attached to the top is invariant under any rotation of the coordinate system.

The problem of determining the wave functions of the stationary states of a symmetric top thus reduces to finding the common eigenfunctions of $\hat{\mathbf J}^{2},\hat J_z,\hat J_\zeta$. This question, in turn, is closely connected mathematically with the transformation law of angular-momentum eigenfunctions under

\SourcePage{498}{501}
finite rotations. Changing the notation for the quantum numbers, we write this law (58.7) as
\begin{equation}
 \psi_{JM}=\sum_k D_{kM}^{(J)}(\alpha,\beta,\gamma)\psi_{Jk}. \tag{103.7}
\end{equation}
We shall take $\psi_{JM}$ to be the wave function of a state of the top described relative to the fixed coordinate axes $xyz$, and $\psi_{Jk}$ to be wave functions of states described relative to the axes $\xi\eta\zeta$ attached to the top. In coordinates rigidly attached to the physical system (the top), however, the quantities $\psi_{Jk}$ have definite values independent of the system's orientation in space; denote them by $\psi_{Jk}^{(0)}$. Formula (103.7) then gives the angular dependence of the functions $\psi_{JM}$. Now let the state $|JM\rangle$ also have a definite value $k$ of the angular-momentum projection on the $\zeta$ axis. This means that only one of all the quantities $\psi_{Jk}^{(0)}$ is nonzero, namely the one with the specified value of $k$. The sum in (103.7) then reduces to one term:
\[
 \psi_{JMk}=\psi_{Jk}^{(0)}D_{kM}^{(J)}(\alpha,\beta,\gamma).
\]

We have thereby found how the wave functions of the states $|JMk\rangle$ depend on the Euler angles specifying the rotation of the top's axes relative to the fixed axes. Normalizing the wave function by
\[
 \int|\psi_{JMk}|^2\sin\beta\,d\alpha\,d\beta\,d\gamma=1,
\]
we obtain
\begin{equation}
 \psi_{JMk}=i^J\sqrt{\frac{2J+1}{8\pi^2}}D_{kM}^{(J)}(\alpha,\beta,\gamma); \tag{103.8}
\end{equation}
the phase factor has been chosen so that, for $k=0$, (103.8) becomes the eigenfunction of a free integer angular momentum $J$ (not tied in any way to the $\zeta$ axis) with projection $M$, that is, the usual spherical function (cf. (58.25)\footnote{For a direct derivation of (103.8), without recourse to the theory of finite rotations, see Problem~1 of this section. For the calculation of matrix elements of various quantities using the wave functions (103.8), see \S~110, 87 (the corresponding formulas differ from those for a diatomic molecule without spin only in the notation of the quantum numbers; cf. the note on p.~389).}).

\SourcePage{499}{502}
\textbf{Asymmetric top.} When $I_A\ne I_B\ne I_C$, the energy levels cannot be calculated in general form. The degeneracy in the directions of the angular momentum relative to the top is completely removed, so that a given $J$ corresponds to $2J+1$ distinct nondegenerate levels. To calculate these levels (for a specified $J$), one starts from the Schrödinger equation written in matrix form (O.~Klein, 1929). This is done as follows.

The wave functions $\psi_{Jk}$ of the states of the top with definite $J$ and definite $\zeta$ projection of the angular momentum are the functions (103.8) found above (the index $M$ for the $z$ projection, on which the energy does not depend, will henceforth be omitted for brevity); in these states the energy of an asymmetric top has no definite value. Conversely, in stationary states the projection $J_\zeta$ has no definite value, so definite values of $k$ cannot be assigned to the energy levels. We seek the wave functions of these states as linear combinations
\begin{equation}
 \psi_J=\sum_k c_k\psi_{Jk} \tag{103.9}
\end{equation}
(all the functions are understood to have some one common value of $M$). Substitution in the Schrödinger equation $\hat H\psi_J=E_J\psi_J$ gives the system
\begin{equation}
 \sum_{k'}\bigl(\langle Jk|H|Jk'\rangle-E\delta_{kk'}\bigr)c_{k'}=0, \tag{103.10}
\end{equation}
and the condition for this system to have a solution yields the secular equation
\begin{equation}
 \left|\langle Jk|H|Jk'\rangle-E\delta_{kk'}\right|=0. \tag{103.11}
\end{equation}
The roots of this equation determine the energy levels of the top. The system (103.10) then gives the linear combinations (103.9) that diagonalize the Hamiltonian, that is, the wave functions of the stationary states of the top with specified $J$ (and $M$). The calculation of a matrix element of any physical quantity from these wave functions is thus reduced to matrix elements for the symmetric top.

The operators $\hat J_\xi,\hat J_\eta$ have matrix elements only for transitions in which $k$ changes by one, while $\hat J_\zeta$ has diagonal elements only (see formulas (27.13), in which $J,k$ are to be written in place of $L,M$). Hence $\hat J_\xi^2,\hat J_\eta^2,\hat J_\zeta^2$, and consequently $\hat H$, have

\SourcePage{500}{503}
matrix elements only for transitions $k\to k,k\pm2$. The absence of matrix elements for transitions between states with even and odd $k$ makes the secular equation of degree $2J+1$ split at once into two independent equations of degrees $J$ and $J+1$. One is formed from the matrix elements for transitions among states with even $k$, and the other from those among states with odd $k$.

Each of these equations can in turn be reduced to two equations of lower degree. For this purpose the matrix elements must be defined not with the functions $\psi_{Jk}$ but with the functions
\begin{equation}
 \begin{aligned}
 \psi_{Jk}^{+}&=\frac1{\sqrt2}(\psi_{Jk}+\psi_{J,-k}),\\
 \psi_{Jk}^{-}&=\frac1{\sqrt2}(\psi_{Jk}-\psi_{J,-k})\quad(k\ne0),\\
 \psi_{J0}^{+}&=\psi_{J0}.
 \end{aligned} \tag{103.12}
\end{equation}
Functions distinguished by the indices $+$ and $-$ have different symmetries (under reflection in a plane through the $\zeta$ axis, which changes the sign of $k$), so matrix elements for transitions between them vanish. Secular equations may therefore be formed separately for the $+$ and $-$ states.

The Hamiltonian (103.1), together with the commutation rules (103.3), has a specific symmetry: it is invariant under a simultaneous change of sign of any two of the operators $\hat J_\xi,\hat J_\eta,\hat J_\zeta$. Formally, this symmetry corresponds to the group $D_2$. The levels of the asymmetric top can therefore be classified by the irreducible representations of this group. There are thus four types of nondegenerate level, corresponding to the representations $A,B_1,B_2,B_3$ (see Table~7, p.~460).

It is easy to establish which states of the asymmetric top belong to each of these types. To do this, the symmetry properties of $\psi_{Jk}$ and of the functions (103.12) constructed from them must be found. They could be obtained directly from (103.8). It is simpler, however, to proceed from the more usual spherical functions, noting that the symmetry properties of the wave functions of states with definite angular-momentum projection on the $\zeta$ axis are the same as those of the angular-momentum eigenfunctions
\begin{equation}
 \psi_{Jk}\sim Y_{Jk}^*(\theta,\varphi)\sim e^{-ik\varphi}\Theta_{Jk}(\theta), \tag{103.13}
\end{equation}

\SourcePage{501}{504}
where $\theta,\varphi$ are the spherical angles in the $\xi\eta\zeta$ axes, and here the sign $\sim$ means ``transforms as''; the complex conjugation in (103.13) is connected with the reversed sign on the right-hand sides of the commutation relations (103.3).

A rotation through $\pi$ about the $\zeta$ axis (that is, the symmetry operation $C_2^{(\zeta)}$) multiplies (103.13) by $(-1)^k$:
\[
 C_2^{(\zeta)}:\qquad \psi_{Jk}\to(-1)^k\psi_{Jk}.
\]
The operation $C_2^{(\eta)}$ may be regarded as the result of an inversion followed by a reflection in the $\xi\zeta$ plane; the first operation multiplies $\psi_{Jk}$ by $(-1)^J$, while the second (a change in the sign of $\varphi$) is equivalent to changing the sign of $k$. Taking account of the definition of $\Theta_{J,-k}$ in (28.6), we therefore obtain
\[
 C_2^{(\eta)}:\qquad \psi_{Jk}\to(-1)^{J+k}\psi_{J,-k}.
\]
Finally, under $C_2^{(\xi)}=C_2^{(\eta)}C_2^{(\zeta)}$ we have
\[
 C_2^{(\xi)}:\qquad \psi_{Jk}\to(-1)^J\psi_{J,-k}.
\]

Using these transformation laws, we find that the states corresponding to (103.12) belong to the following symmetry types:
\begin{equation}
\begin{array}{c@{\quad}c@{\quad}c}
\psi_{Jk}^{+}\left\{\begin{array}{ll}
\text{even }J,\ \text{even }k&-A,\\
\text{even }J,\ \text{odd }k&-B_3,\\
\text{odd }J,\ \text{even }k&-B_1,\\
\text{odd }J,\ \text{odd }k&-B_2,
\end{array}\right.\\[6pt]
\psi_{Jk}^{-}\left\{\begin{array}{ll}
\text{even }J,\ \text{even }k&-B_1,\\
\text{even }J,\ \text{odd }k&-B_2,\\
\text{odd }J,\ \text{even }k&-A,\\
\text{odd }J,\ \text{odd }k&-B_3.
\end{array}\right.
\end{array} \tag{103.14}
\end{equation}
A simple count gives the number of states of each type for a specified $J$. In particular, the following numbers of states correspond to type $A$ and to each of the types $B_1,B_2,B_3$:
\begin{equation}
\begin{array}{c|c|c}
&A&B_1,B_2,B_3\\ \hline
\text{Even }J&J/2+1&J/2\\
\text{Odd }J&(J-1)/2&(J+1)/2
\end{array} \tag{103.15}
\end{equation}

\SourcePage{502}{505}
For an asymmetric top there are selection rules for matrix elements governing transitions between states of the types $A,B_1,B_2,B_3$; these follow in the usual way from symmetry considerations. Thus, for the components of a vector physical quantity $\mathbf A$, the selection rules are
\begin{equation}
\begin{aligned}
\text{for }A_\xi:&\quad A\leftrightarrow B_3^{(\xi)},\quad B_1^{(\zeta)}\leftrightarrow B_2^{(\eta)},\\
\text{for }A_\eta:&\quad A\leftrightarrow B_2^{(\eta)},\quad B_1^{(\zeta)}\leftrightarrow B_3^{(\xi)},\\
\text{for }A_\zeta:&\quad A\leftrightarrow B_1^{(\zeta)},\quad B_2^{(\eta)}\leftrightarrow B_3^{(\xi)}
\end{aligned} \tag{103.16}
\end{equation}
(for clarity, the axis about which a rotation has character $+1$ in the given representation is indicated as an index on the representation symbol).

\subsection*{Problems}

\ProblemHeading{1.} Find the wave functions of the states $|JMk\rangle$ of a symmetric top by direct calculation as common eigenfunctions of $\hat{\mathbf J}^{2},\hat J_z,\hat J_\zeta$ (F.~Reiche, H.~Rademacher, 1926).

\SolutionHeading We wish to obtain $\psi_{JMk}$ as a function of the Euler angles $\alpha,\beta,\gamma$ and must therefore express in terms of them the operators for the projections of angular momentum on the fixed axes $x,y,z$. Since the operator for the projection of angular momentum on any axis is $-i\partial/\partial\varphi$, where $\varphi$ is the angle of rotation about that axis, we may write
\[
 \hat J_x=-i\frac{\partial}{\partial\varphi_x},\qquad
 \hat J_y=-i\frac{\partial}{\partial\varphi_y},\qquad
 \hat J_z=-i\frac{\partial}{\partial\varphi_z},
\]
where $\varphi_x,\varphi_y,\varphi_z$ are the angles of rotation about the corresponding axes. Derivatives with respect to these angles can be expressed in terms of derivatives with respect to $\alpha,\beta,\gamma$ by recalling that infinitesimal rotations add as vectors directed along the axes of rotation. The directions of the infinitesimal-rotation vectors $\delta\alpha,\delta\beta,\delta\gamma$ described by the Euler angles are shown in Fig.~20 (p.~270). Projecting them on the fixed axes $xyz$, we find the angles of rotation about these axes in the form
\[
\begin{aligned}
 \delta\varphi_x&=-\sin\alpha\,\delta\beta+\cos\alpha\sin\beta\,\delta\gamma,\\
 \delta\varphi_y&= \cos\alpha\,\delta\beta+\sin\alpha\sin\beta\,\delta\gamma,\\
 \delta\varphi_z&=\delta\alpha+\cos\beta\,\delta\gamma.
\end{aligned}
\]
Conversely,
\[
\begin{aligned}
 \delta\alpha&=-\cot\beta\cos\alpha\,\delta\varphi_x-\cot\beta\sin\alpha\,\delta\varphi_y+\delta\varphi_z,\\
 \delta\beta&=-\sin\alpha\,\delta\varphi_x+\cos\alpha\,\delta\varphi_y,\\
 \delta\gamma&=\frac{\cos\alpha}{\sin\beta}\delta\varphi_x+\frac{\sin\alpha}{\sin\beta}\delta\varphi_y.
\end{aligned}
\]
These expressions give
\[
\begin{aligned}
 \hat J_x&=-i\left(-\cos\alpha\cot\beta\frac{\partial}{\partial\alpha}-\sin\alpha\frac{\partial}{\partial\beta}+\frac{\cos\alpha}{\sin\beta}\frac{\partial}{\partial\gamma}\right),\\
 \hat J_y&=-i\left(-\sin\alpha\cot\beta\frac{\partial}{\partial\alpha}+\cos\alpha\frac{\partial}{\partial\beta}+\frac{\sin\alpha}{\sin\beta}\frac{\partial}{\partial\gamma}\right),\\
 \hat J_z&=-i\frac{\partial}{\partial\alpha}.
\end{aligned}
\]

\SourcePage{503}{506}
When acting on $\psi_{JMk}$, the operators $\hat J_z=-i\partial/\partial\alpha$ and $\hat J_\zeta=-i\partial/\partial\gamma$ ($\gamma$ is the angle of rotation about the $\zeta$ axis!) are replaced by $M$ and $k$ (the corresponding dependence of the wave function on the Euler angles $\alpha$ and $\gamma$ is given by the factor $\exp(i\alpha M+i\gamma k)$). We then have
\[
\begin{aligned}
 \hat J_+&=\hat J_x+i\hat J_y=e^{i\alpha}\left(\frac{\partial}{\partial\beta}-M\cot\beta+\frac{k}{\sin\beta}\right),\\
 \hat J_-&=\hat J_x-i\hat J_y=e^{-i\alpha}\left(-\frac{\partial}{\partial\beta}-M\cot\beta+\frac{k}{\sin\beta}\right).
\end{aligned}
\]
The rest of the derivation is exactly analogous to that given at the end of \S~28. We begin with $\hat J_+\psi_{JJk}=0$, which holds for the wave function with $M=J$. Hence
\[
 \left(\frac{\partial}{\partial\beta}-J\cot\beta+\frac{k}{\sin\beta}\right)\psi_{JJk}=0.
\]
The normalized solution of this equation is
\[
 \psi_{JJk}=i^J(-1)^{J-k}\left[\frac{(2J+1)!}{2(J+k)!(J-k)!}\right]^{1/2}
 \left(\cos\frac\beta2\right)^{J+k}\left(\sin\frac\beta2\right)^{J-k}
 \frac{e^{i(J\alpha+k\gamma)}}{2\pi}
\]
(the normalization integral reduces to Euler's $B$ integral). This expression does indeed coincide, up to a phase factor, with
\[
 \sqrt{\frac{2J+1}{8\pi^2}}D_{kJ}^{(J)}(\alpha,\beta,\gamma)
\]
(cf. (58.26)); the phase factor has been chosen in accordance with the definition in (103.7).

The wave functions with $M<J$ are then calculated by repeated application to $\psi_{JJk}$ of the formula
\[
 \hat J_-\psi_{J,M+1,k}=\sqrt{(J-M)(J+M+1)}\psi_{JMk}.
\]
The final result coincides with (103.8), where the functions $D_{kM}^{(J)}$ are given by (58.9)--(58.11) (with account taken of their symmetry property (58.18)).

\ProblemHeading{2.} Calculate the matrix elements $\langle Jk'|H|Jk\rangle$ for an asymmetric top.

\SolutionHeading From formulas (27.13) we find
\[
 \langle k|J_\xi^2|k\rangle=\langle k|J_\eta^2|k\rangle=\frac12[J(J+1)-k^2],
\]
\[
\begin{aligned}
 \langle k|J_\xi^2|k+2\rangle=\langle k+2|J_\xi^2|k\rangle
 &=-\langle k|J_\eta^2|k+2\rangle=-\langle k+2|J_\eta^2|k\rangle\\
 &=\frac14\sqrt{(J-k)(J-k-1)(J+k+1)(J+k+2)}
\end{aligned}
\]
(the diagonal indices $J,J$ on the matrix elements are omitted throughout for brevity). Hence the required matrix elements of the

\SourcePage{504}{507}
Hamiltonian are\footnote{In Problems 2--5, to simplify the notation in the formulas, we use $a=1/I_A$, $b=1/I_B$, $c=1/I_C$.}
\begin{align}
 \langle k|H|k\rangle&=\frac{\hbar^2}{4}(a+b)[J(J+1)-k^2]+\frac{\hbar^2}{2}ck^2,\notag\\
 \langle k|H|k+2\rangle=\langle k+2|H|k\rangle&=\frac{\hbar^2}{8}(a-b)\sqrt{(J-k)(J-k-1)(J+k+1)(J+k+2)}. \tag{1}
\end{align}
The matrix elements with respect to the functions (103.12) are expressed in terms of the elements (1) by
\begin{align}
 \langle k\pm|H|k\pm\rangle&=\langle k|H|k\rangle,\quad k\ne1,\notag\\
 \langle1\pm|H|1\pm\rangle&=\langle1|H|1\rangle\pm\langle1|H|-1\rangle,\notag\\
 \langle k\pm|H|k+2,\pm\rangle&=\langle k|H|k+2\rangle,\quad k\ne0,\notag\\
 \langle0+|H|2+\rangle&=\sqrt2\langle0|H|2\rangle. \tag{2}
\end{align}

\ProblemHeading{3.} Determine the energy levels of an asymmetric top for $J=1$.

\SolutionHeading The secular equation of third degree splits into three linear equations. One of them gives
\begin{equation}
 E_1=\langle0+|H|0+\rangle=\frac{\hbar^2}{2}(a+b). \tag{3}
\end{equation}
The other two energy levels may be written down immediately, since it is evident in advance that the three parameters $a,b,c$ enter the problem symmetrically. Thus
\begin{equation}
 E_2=\frac{\hbar^2}{2}(a+c),\qquad E_3=\frac{\hbar^2}{2}(b+c). \tag{4}
\end{equation}
The levels $E_1,E_2,E_3$ belong respectively to the symmetry types $B_1,B_2,B_3$\footnote{This follows directly from symmetry considerations. For example, the energy $E_1$ is symmetric in the parameters $a$ and $b$; this is required for the energy of a state having the same symmetry with respect to the $\xi$ and $\eta$ axes (a state of type $B_1$).}. The wave functions of these states are
\[
 \psi_1=\psi_{10}^+,\qquad \psi_2=\psi_{11}^-,\qquad \psi_3=\psi_{11}^+.
\]

\ProblemHeading{4.} The same for $J=2$.

\SolutionHeading The secular equation of fifth degree splits into three linear equations and one quadratic equation. One of the linear equations gives
\begin{equation}
 E_1=\langle2-|H|2-\rangle=2\hbar^2c+\frac{\hbar^2}{2}(a+b) \tag{5}
\end{equation}
(a level of type $B_1$). We infer from this that there must be two further levels (of types $B_2$ and $B_3$):
\[
 E_2=2\hbar^2b+\frac{\hbar^2}{2}(a+c),\qquad E_3=2\hbar^2a+\frac{\hbar^2}{2}(b+c).
\]

\SourcePage{505}{508}
The wave functions corresponding to these three levels are
\[
 \psi_1=\psi_{22}^-,\qquad \psi_2=\psi_{21}^-,\qquad \psi_3=\psi_{21}^+.
\]
The quadratic equation is
\begin{equation}
 \begin{vmatrix}
 \langle0+|H|0+\rangle-E&\langle2+|H|0+\rangle\\
 \langle2+|H|0+\rangle&\langle2+|H|2+\rangle-E
 \end{vmatrix}=0. \tag{6}
\end{equation}
Solving it, we obtain
\begin{equation}
 E_{4,5}=\hbar^2(a+b+c)\pm\hbar^2[(a+b+c)^2-3(ab+bc+ac)]^{1/2}. \tag{7}
\end{equation}
These levels are of type $A$. Their wave functions are linear combinations of $\psi_{20}^+$ and $\psi_{22}^+$.

\ProblemHeading{5.} The same for $J=3$.

\SolutionHeading The secular equation of seventh degree splits into one linear and three quadratic equations. The linear equation gives
\begin{equation}
 E_1=\langle2-|H|2-\rangle=2\hbar^2(a+b+c) \tag{8}
\end{equation}
(a level of type $A$). One of the quadratic equations is equation (6) of the preceding problem (with a different value of $J$). Its roots are
\begin{equation}
 E_{2,3}=\frac{5\hbar^2}{2}(a+b)+\hbar^2c\pm\hbar^2[4(a-b)^2+c^2+ab-ac-bc]^{1/2} \tag{9}
\end{equation}
(levels of type $B_1$). The remaining levels are obtained from these by permuting the parameters $a,b,c$.

\ProblemHeading{6.} Determine the splitting of the levels of a system possessing a quadrupole moment in an arbitrary external electric field.

\SolutionHeading Choose the principal axes of the tensor $\partial^2\varphi/\partial x_i\partial x_k$ as coordinate axes (cf. Problem~3, \S~76). The quadrupole part of the system's Hamiltonian is then reduced to the form
\[
 \hat H=A\hat J_x^2+B\hat J_y^2+C\hat J_z^2,\qquad A+B+C=0.
\]
Because this expression is completely analogous in form to the Hamiltonian (103.1), the stated problem is equivalent to finding the energy levels of an asymmetric top. The only differences are that the sum of the coefficients is now $A+B+C=0$ and that the angular momentum may also take half-integer values. For the latter, the calculation must be carried out anew by the same method, while for integer $J$ the results of Problems 3--5 may be used. We finally obtain the following energy shifts $\Delta E$ for the first few values of $J$:
\[
\begin{array}{c@{\qquad}l}
J=1:&\Delta E=-A,-B,-C,\\
J=3/2:&\Delta E=\pm\sqrt{(3/2)(A^2+B^2+C^2)},\\
J=2:&\Delta E=3A,3B,3C,\ \pm\sqrt{6(A^2+B^2+C^2)}.
\end{array}
\]
For $J=3/2$ the energy levels remain twofold degenerate, in accordance with Kramers' theorem (\S~60).
